%% !BIB TS-program = biber
\documentclass[11pt]{amsart}
\usepackage[a4paper]{geometry}
\usepackage[backend=biber,autocite=plain,backref=true]{biblatex}
\usepackage{amsmath, amssymb, amsfonts}
\usepackage{url}
\addbibresource{conical_shell.bib}

\author{Maximillian Julius Otto Strutt \cite{strutt_1933}}
\title{Natural Frequencies of a Thin Conical Shell}
\address[Eindhoven]{Natuurkundig Laboratorium der N. V. Philips' Gloeilampenfabrieken, Eindhoven/Holland}
\date{May 18 1933}
\begin{document}
\maketitle

\textbf{Summary}
The natural frequencies of a thin conical shell are calculated under the
assumption that no vibrations cause expansion of the shell.  The calculation
method is from Lord Rayleigh \cite[\S 235 f]{rayleigh_sound_v1_1945},
\cite[page 551]{rayleigh_1902}, \cite{rayleigh_1881}.  To approximate the
conditions practically encountered in loudspeaker cones, it is assumed that the
conical shell is enclosed on one side by a rigid disk and on the other side
free.  The transition from the cone to the known formulas for the cylinder is
carried out.  A simple expression for the natural frequencies is obtained asymptotically.

\setcounter{page}{729}

\section{Introduction}
The moving part (diaphragm) of loudspeakers is often designed in the shape of a
cone for reasons of stiffness.

The differential equations for the deformation of a conical
shell~\cite[page~603]{love_edition4} are so complex that no attempt has been
made to calculate the natural frequencies based on these equations.  In
contrast, the calculation using Rayleigh's method~\cite{rayleigh_1902} for a
thin shell (which applies to loudspeakers) is straightforward.

If coordinates are introduced on the surface (\(x\) in
the generating direction towards the apex), \(y\) in the
tangent direction perpendicular to the axis, \(\bar{z}\)
in the direction perpendicular to the generating
direction and inwards to \(y\), then before deformation:
\[
\bar{z}= \frac{1}{2}
\left(\frac{x^2}{\rho_1} + \frac{y^2}{\rho_2} \right).
\]
with \(\rho_1\) and \(\rho_2\) the principal radii of
curvature (\(\rho_1 = \infty\)).

After deformation, however:
\[
\bar{z} = \frac{1}{2} \left(
\frac{x^2}{\rho_1 + d\rho_1} +
\frac{y^2}{\rho_2 + d\rho_2} + 2\tau x y
\right).
\]

Suppose that \(2d\) is the shell thickness, \(E\) is the
modulus of elasticity, \(\sigma\) is the Poisson constant.
\[
E = \frac{-3n}{m - n/3}, \quad \sigma = \frac{(m - n)}{2m}.
\]
If \(\delta (1/\rho_1)\), \(\delta (1/\rho_2)\), and
\(\tau\) are known, the potential energy per unit area
of the shell is written as:
\begin{equation} \label{eq1}
U = \frac{2n d^2}{3} \left(
\left( \delta \frac{1}{\rho_1} \right)^2 +
\left( \delta \frac{1}{\rho_2} \right)^2 +
\frac{m-n}{m+n} \left(\delta \frac{1}{\rho_1} +
\delta \frac{1}{\rho_2}\right)^2 + 2 \tau^2 \right).
\end{equation}
Integration over the shell area yields the total
potential energy.  If the kinetic energy is also known,
the natural frequencies can be calculated.

\section{Calculation of Potential Energy}

The first step is to express the aforementioned surface
coordinates in cylindrical coordinates, measured from the cone apex, and \(r\)
and \(\varphi\) measured from the considered lateral surface point \( (r_o,z_o, 0)\).

The formulas are (\(\beta \) = angle axis-generating line):
\begin{align}
x &= \frac{r_o}{\sin(\beta)} - r \cos(\beta) \sin(\beta) - z \cos(\beta) \nonumber \\
y &= r \sin(\varphi) \label{eq2} \\
\bar{z} &= 2 \sin(\beta ) - r \cos(\varphi) \cos(\beta) \nonumber
\end{align}

Near to the point \((r_o,z_o, 0)\) (except for
second-order terms in \(\varphi\) and \(z-z_o\))
\begin{align*}
x&= \frac{r-r_o}{\sin(\beta)}=
\frac{z-z_o}{\sin(\beta)} \\
y &= r \varphi \tag{2a} \\
\bar{z} & = 0.
\end{align*}

If \(r\), \(z\), \(\varphi\) each receive a small
increment \(u\), \(v\), \(w\), the following equations result:
\begin{align}
x &= \frac{r_o}{\sin(\beta)} -z_o \cos(\beta)
-r_o\cos(\beta)\sin(\beta) -[u_o\sin(\beta)] \nonumber \\
& -\left(\frac{d u}{dz}\right)_o (z-z_o) \sin(\beta) -
\left(\frac{d u}{d\varphi}\right)_o
\varphi ~~ \sin(\beta) -[v_o\cos(\beta)] \\
& -\left(\frac{d v}{dz}\right)_o (z-z_o)\cos(\beta) -
\left(\frac{d v}{d \varphi}\right)_o
\varphi \cos(\beta) + [r_o w_o \varphi \sin(\beta)]; \nonumber \\
u_o w_o &=(r-r_o)w_o = (z-z_o)\tan(\beta)w_o; \nonumber
\end{align}

Equation 3 continues on the next page.  \setcounter{equation}{2}
\begin{equation}\label{eq3p731}
\begin{aligned}
y & =r_0 \varphi+\left[r_0 w_0\right]+u_0 \varphi+r_0\left(\frac{d w}{d z}\right)_0\left(z-z_0\right)+r_0\left(\frac{d w}{d \varphi}\right)_0 \varphi; \\
\bar{z} & =z_0 \sin \beta-r_0 \cos \beta+\left[v_0 \sin \beta\right]-\left[u_0 \cos \beta\right] \\
& +\left(\frac{d v}{d z}\right)_0\left(z-z_0\right) \sin \beta+\left(\frac{d v}{d \varphi}\right)_0 \varphi \sin \beta \\
& -\left(\frac{d u}{d z}\right)_0\left(z-z_0\right) \cos \beta-\left(\frac{d u}{d \varphi}\right)_0 \varphi \cos \beta+r_0 w_0 \varphi \cos \beta .
\end{aligned}
\end{equation}

Terms of the second and higher order in \(\varphi\) and \(z-z_o\) are neglected
here.  The subscript \(0\) indicates that the relevant quantities are formed at
the point, \(r = r_o\), \(z = z_o\), \(\varphi = 0\).

The surface can now be returned to its original position (neglecting terms of
second and higher order) by means of an infinitesimal translation and rotation.
The translation is equal and opposite to the terms enclosed in square brackets in (3).

The additional rotational terms for \(x\), \(y\), and
\(\bar{z}\) are, respectively:
\[r_o \phi \omega_3;\quad (z - z_o)
\frac{\omega_3}{\cos \beta};\quad -(z - z_o)
\frac{\omega_3}{\cos \beta} - r_o \phi \omega_1\]

A necessary condition for regaining the zero position described above is that,
\begin{subequations}\label{eq4abcdef}
\begin{align}
r \omega_3-\frac{d u}{d \varphi} \sin \beta-\frac{d v}{d \varphi} \cos \beta+r w \sin \beta & =0; \\
\frac{d u}{d z} \tan \beta+\frac{d v}{d z} & =0; \\
u \quad+r \frac{d w}{d \varphi} & =0; \\
w \tan \beta + r \frac{d w}{d z}+\frac{\omega_3}{\cos \beta}=0; \\
r w \cos \beta+\frac{d v}{d \varphi} \sin \beta-\frac{d u}{d \varphi} \cos \beta-r \omega_1 & =0; \\
-\frac{d v}{d z} \sin \beta+\frac{d u}{d z} \cos \beta+\frac{\omega_2}{\cos \beta}=0 . &
\end{align}
\end{subequations}
(Index 0 omitted) Eliminating \(\omega_3\)  from (4d) and (4a) yields:
\[
\tag{4ad} \qquad
r^2 \frac{d w}{d z}+\frac{d u}{d q} \tan \beta+\frac{d v}{d \varphi}=0
\]
It is easy to verify that equations (4b), (4c), and (4ad) are identical to
those obtained by Lord Rayleigh \cite[p 399 Eq 14]{rayleigh_sound_v1_1945}
by a different method.
Their solution, under the condition that the cone's surface is enclosed by an
inextensible disk at \(z = l_1\), \(u(x,y, l_1)=0\) and \(w(x,y,l_1)=0\),
is as follows: For \( i \geq 2 \),
\begin{equation}\label{eq5p732}
\begin{aligned}
u=&\sum_i i \tan \beta\left(z-l_1\right) A_i \sin i \varphi; \\
v=&\sum_i-\tan^2 \beta\left( \frac{l_1}{i}+i\left(z-l_1\right)\right) A_i \sin i \varphi, \\
w=&\sum_i\left(1-\frac{l_1}{z}\right) A_i \cos i \varphi .
\end{aligned}
\end{equation}

To calculate the potential energy according to Equation (1), it is necessary to take
into account all second-order terms in $\phi$ and $z - z_0$ in the final equation
(3) (after performing the infinitesimal rotation described above).

Taking \(\omega_1\) and \(\omega_2\) from Equations (4e) and (4f) respectively,
\[
\begin{aligned}
2 \bar{z} & =\left\{\frac{d^2 v}{d z^2}\left(z-z_0\right)^2+\frac{d^2 v}{d \varphi^2} \varphi^2+2 \frac{d^2 v}{d \varphi}\left(z-z_0\right) \varphi\right\} \sin \beta \\
& -\left\{-u \varphi^2-r \varphi^2+\frac{d^2 u}{d z^2}\left(z-z_0\right)^2+2 \frac{d^2 u}{d z d \varphi}\left(z-z_0\right) \varphi\right. \\
& \left.+\frac{d^2 u}{d \varphi^2} \varphi^2-2 r \frac{d w}{d z}\left(z-z_0\right) \varphi-2 r \frac{d w}{d \varphi} \varphi^2\right\} \cos \beta \\
& +2 \frac{d v}{d z} \varphi^2 \sin ^2 \beta \cos \beta-2 \frac{d u}{d z} \varphi^2 \sin \beta \cos ^2 \beta .
\end{aligned}
\]
Using \(x = (z-z_o)/ \cos \beta\), \(y = r \phi\), it follows from this that:
\[
\begin{aligned}
2\bar{z}= &-2xy \frac{\cos \beta}{r}\left\{\frac{d^2 v}{d z d \varphi} \sin \beta-\frac{d^2 u}{d z d \varphi} \cos \beta+r \frac{d w}{d z} \cos \beta\right\} \\
& +\frac{y^2}{r^2}\left\{\frac{d^2 v}{d \varphi^2} \sin \beta+u \cos \beta+r \cos \beta-\frac{d^2 u}{d \varphi^2} \cos \beta\right. \\
& \left.+2 r \frac{d w}{d \varphi} \cos \beta-r \frac{d u}{d z} \sin \beta \cos ^2 \beta+r \frac{d v}{d z} \sin ^2 \beta \cos \beta\right\} \\
& +x^2 \cos ^2 \beta\left\{\frac{d^2 v}{d z^2} \sin \beta-\frac{d^2 u}{d z^2} \cos \beta\right\} .
\end{aligned}
\]
Taking into account Equations (4b), (4c), and (4ad), one can simplify to:
\begin{equation}\label{eq6p732}
\left\{\begin{aligned}
2 \bar{z}= & -2 x y \frac{\cos \beta}{r}\left\{\frac{d^2 v}{d z d \varphi} \sin \beta-\frac{d^2 u}{d z d \varphi} \cos \beta+r \frac{d w}{d z} \cos \beta\right\} \\
& +\frac{y^2}{r^2}\left\{\frac{d^2 v}{d \varphi^2} \sin \beta+r \cos \beta-u \cos \beta-\frac{d^2 u}{d \varphi^2} \cos \beta\right. \\
& \left.-r \frac{d u}{d z} \sin \beta \cos ^2 \beta+r \frac{d v}{d z} \sin ^2 \beta \cos \beta\right\} .
\end{aligned}\right.
\end{equation}
The transition to the cylinder \(beta = 0\) leads to
Lord Rayleigh's formulas \cite[p 414 Eq 6]{rayleigh_sound_v1_1945}.

Before deformation, (6) read:
\[2\bar{z}=y^2 \frac{\cos \beta}{r^2}.\]
The quantities
\(\delta \frac{1}{e_1}\), \(\delta \frac{1}{e_2} \)
and \(r\) in Equation (1) are thus:
\begin{equation}\label{eq7p733}
\begin{aligned}
\delta\left(\frac{1}{\rho_1}\right)= & 0; \\
\delta\left(\frac{1}{\rho_2}\right)= & \frac{1}{r^2} \frac{d^2 v}{d \varphi^2} \sin \beta-\frac{d^2 u}{d \varphi^2} \frac{\cos \beta}{r^2}-\frac{u \cos \beta}{r^3} \\
& -\frac{\sin \beta \cos ^2 \beta}{r} \frac{d u}{d z}+\frac{\sin ^2 \beta \cos \beta}{r} \frac{d v}{d z}; \\
-\tau= & \frac{\cos \beta}{r}\left\{\frac{d^2 v}{d z d \varphi} \sin \beta-\frac{d^2 u}{d z d \varphi} \cos \beta
+r \frac{d w}{d z} \cos \beta\right\} .
\end{aligned}
\end{equation}


Substituting the solutions (5) into (7) and the resulting expressions into (1)
after integration over \( 0 \leq \varphi \leq 2 \pi \) yields:
\begin{equation}\label{eq8p733}
\begin{aligned}
\delta\left(\frac{1}{\varrho_2}\right) & =\sum_i A_i \sin(i \varphi) ~ \alpha_i\left(\frac{1}{z}-\frac{l_1}{z^2}\right); \\
\alpha_i & =\frac{i^3-i}{\sin \beta}; \\
-\tau & =\sum_i A_i \cos( i \varphi) ~ \frac{1}{z} \cdot\left(\frac{l_1}{z} \cos ^2 \beta-i^2\right); \\
U & =\frac{4 \pi n d^2}{3} \sum_i A_i^2\left\{\frac{m}{m+n} \alpha_i^2\left(\frac{1}{z}-\frac{l_1}{z^2}\right)^2
 +\frac{1}{z^2}\left(\frac{l_1}{z} \cos ^2 \beta-i^2\right)\right\} .
\end{aligned}
\end{equation}
From this, the total potential energy is obtained by integration over \(z\):
\begin{equation}\label{eq9p733}
\begin{aligned}
V= & \int_{l_1}^{l_2} d z \cdot z \tan \beta \cdot U \\
& =\frac{4 \pi n d^3}{3}i \tan \beta \sum_i A_i^2\left[\frac { m \alpha _ { i } ^ { 2 } } { m + n } \left\{\ln \frac{l_2}{l_1}\right.\right. \\
& \left.+2\left(\frac{l_1}{l_2}-1\right)-\frac{1}{2}\left(\frac{l_1^2}{l_2^2}-1\right)\right\}-\left(\frac{l_1^2}{l_2^2}-1\right) \frac{\cos ^4 \beta}{2} \\
& \left.+2 i^2 \cos ^2 \beta\left(\frac{l_1}{l_2}-1\right)+i^4 \ln \frac{l_2}{l_2}\right] .
\end{aligned}
\end{equation}
This expression applies to static problems.

\section{Natural Frequencies}
The cone's surface has kinetic energy,
\[
T=\varrho d \int_{l_1}^{l_2} d z ~ z \tan \beta \int_0^{2 \pi} d \varphi\left\{\left(\frac{d u}{d t}\right)^2+\left(\frac{d v}{d t}\right)^2+\left(\frac{r d w}{d t}\right)^2\right\}
\]
where \(\rho\) is the density.  Assuming temporal variaiont of \(u,v,w\) of \(\cos \omega t\):
\begin{equation}\label{eq10p734}
\begin{aligned}
-T & = \pi \omega^2 \varrho ~ d ~ l_1^4 \tan^3 \beta \cdot \pi \sum_i A_i^2\left[( i ^ { 2 } + 1 ) \left\{\frac{1}{4}\left(\frac{l_2^4}{l_1^4}-1\right)\right.\right. \\
& \left.+\frac{1}{2}\left(\frac{l_2^2}{l_1^2}-1\right)-\frac{2}{3}\left(\frac{l_2^3}{l_1^3}-1\right)\right\}+\tan^2 \beta\left\{\frac{1}{2 i^2}\left(\frac{l_2^2}{l_1^2}-1\right)\right. \\
& +\frac{i^2}{4}\left(\frac{l_2^4}{l_1^4}-1\right)+\frac{i^2}{2}\left(\frac{l_2^2}{l_1^2}-1\right)+\frac{2}{3}\left(\frac{l_2^3}{l_1^3}-1\right) \\
& \left.\left.-\left(\frac{l_2^2}{l_1^2}-1\right)-\frac{2 i^2}{3}\left(\frac{l_2^3}{l_1^3}-1\right)\right\}\right] .
\end{aligned}
\end{equation}

From equations (9) and (10), the desired angular frequencies
( \(\omega=2 \pi\) Hz) are immediately obtained:
\begin{equation}\label{eq11p734}
\omega_i^2=\frac{4 n d^2}{3 \varrho l_i^4 \tan^2 \beta} \cdot \frac{F_i}{G_i}:
\end{equation}
\begin{equation}\label{eq12p734}
\begin{aligned}
F_i & =\frac{m \alpha_i^2}{m+n}\left\{\ln \frac{1}{x}+2(x-1)-\frac{1}{2}\left(x^2-1\right)\right\} \\
& -\frac{\cos ^4 \beta}{2}\left(x^2-1\right)+2 i^2 \cos ^2 \beta\left(x^2-1\right)+i^4 \ln \frac{1}{x} ;
\end{aligned}
\end{equation}
\begin{equation}\label{eq13p734}
\begin{aligned}
G_i & =\left(i^2+1\right)\left\{\frac{1}{4}\left(\frac{1}{x^4}-1\right)+\frac{1}{2}\left(\frac{1}{x^2}-1\right)\right. \\
& -\frac{2}{3}\left(\frac{1}{x^3}-1\right)+\tan^2 \beta\left\{\frac{1}{2 i^2}\left(\frac{1}{x^2}-1\right)+\frac{i^2}{4}\left(\frac{1}{x^4}-1\right)\right. \\
& +\frac{i^2}{2}\left(\frac{1}{x^2}-1\right)+\frac{2}{3}\left(\frac{1}{x^3}-1\right)-\left(\frac{1}{x^2}-1\right) \\
& \left.-\frac{2 i^2}{3}\left(\frac{1}{x^3}-1\right)\right\}
\end{aligned}
\end{equation}
From (8), for \(i\geq 2\), \(\alpha_i=\frac{i^3-i}{\sin \beta}\), and \(x=l_1/l_2\).

In the case of a cylinder of radius \(R\), \(l_1 \tan \beta \sim l_1 \sin
\beta=R\).  This results in:
\begin{equation}
\omega_i^2=\frac{4 m n d^2}{3 \rho R^4(m+n)}
\left(\frac{\left(i^3-i\right)^2}{i^2+1}\right) \lim_{x \rightarrow 1} \frac{\ln \frac{1}{x}+2(x-1)-\frac{1}{2}\left(x^2-1\right)}{\frac{1}{4}\left(\frac{1}{x^4}-1\right)+\frac{1}{2}\left(\frac{1}{x^2}-1\right)-\frac{2}{3}\left(\frac{1}{x^3}-1\right)},
\end{equation}
or,
\[
\omega_i^2=\frac{4 d^2}{3 \varrho R^4} \frac{m n}{m+n} \frac{\left(i^3-i\right)^2}{i^2+1},
\]
a well-known formula \cite[p 417 Eq 18]{rayleigh_sound_v1_1945}.

It is often useful to have a simple formula at hand for high frequencies. From
(11), (12), (13), and (8), we obtain:
\[
\lim_{i \rightarrow \infty} \omega_i^2=\frac{4 d^2}{3 \varrho l_i^4 \tan^4 \beta} \frac{m n}{m+n}
i^4 \frac{\ln \frac{1}{x}+2(x-1)-\frac{1}{2}\left(x^2-1\right)}{\frac{1}{4}\left(\frac{1}{x^4}-1\right)+\frac{1}{2}\left(\frac{1}{x^2}-1\right)-\frac{2}{3}\left(\frac{1}{x^3}-1\right)}
+ \mathcal{O}(i^2)
\]

I would like to point out that the calculation method used does not allow us to
handle the case of a cone shell extending to the apex.

If the cone is held at \(l_2\) instead of \(l_1\) by a non-extensible disk,
Equation (5) must be simply modified, which also changes (12) and (13)
somewhat. It is left to the reader to handle this and possibly other special
cases using the formulas above.

\printbibliography[heading=bibintoc,title={References}]
\end{document}
